% Written by Claude (2026-11-14): how far the last token can see after three
% layers of sliding-window attention with w = 3 (reach 3(w-1) = 6 tokens).
\begin{tikzpicture}[x=0.95cm, y=0.72cm,
        tok/.style={circle, draw=drawgray!80, fill=white, minimum size=0.42cm, inner sep=0pt},
        hot/.style={tok, draw=drawblue, fill=drawblue_l!80},
        edge/.style={drawblue!70, line width=0.5pt},
        rlab/.style={font=\scriptsize, text=black!60, anchor=east},
    ]
    % Row r (0 = input, 3 = last layer): token i is highlighted if i >= 9 - 2(3 - r).
    \foreach \r/\name in {0/input, 1/layer 1, 2/layer 2, 3/layer 3} {
            \node[rlab] at (-0.5,\r) {\name};
            \foreach \i in {0,...,9} {
                    \pgfmathtruncatemacro{\lo}{9-2*(3-\r)}
                    \ifnum\i<\lo
                        \node[tok] (n\r-\i) at (\i,\r) {};
                    \else
                        \node[hot] (n\r-\i) at (\i,\r) {};
                    \fi
                }
        }
    % Edges from each highlighted token to the w = 3 tokens it attends to.
    \foreach \r in {1,2,3} {
            \pgfmathtruncatemacro{\lo}{9-2*(3-\r)}
            \pgfmathtruncatemacro{\b}{\r-1}
            \foreach \i in {\lo,...,9} \foreach \d in {0,1,2} {
                    \pgfmathtruncatemacro{\j}{\i-\d}
                    \draw[edge] (n\r-\i) -- (n\b-\j);
                }
        }
    \draw[decorate, decoration={brace, amplitude=5pt, mirror}, drawgray]
    (2.8,-0.4) -- (9.2,-0.4) node[midway, below=6pt, font=\scriptsize, text=black!70] {$3(w-1)=6$ tokens in the past};
\end{tikzpicture}
